\section{The intrinsic ample cone and openness on actual finite line families}
\label{sec:native-numerical-ample-cone}

\begin{definition}[Cone of actual geometrically ample line classes {\cite[\S1.2]{BJ22}}]
\label{def:native-numerical-ample-cone}
\lean{EnvelopeProperty.NativeRealNumericalClasses.ampleCone}
\leanok
\uses{thm:native-numerical-ample-classes}
In the original real numerical quotient, define the ample cone to be the
convex cone generated by numerical classes of actual original geometrically
ample invertible lines. Thus the generating geometry is the original
positive-power projective-embedding definition of ampleness.
\end{definition}
\begin{proof}
Take the convex-cone hull of the set of those original line classes. The
cone operations are finite addition and multiplication by strictly positive
real scalars. This definition introduces no numerical positivity premise.
\end{proof}

\begin{lemma}[Actual ample lines and nef containment {\cite[\S1.2]{BJ22}}]
\label{lem:native-numerical-ample-cone-nef}
\lean{EnvelopeProperty.NativeRealNumericalClasses.ampleLineClass_mem_ampleCone, EnvelopeProperty.NativeRealNumericalClasses.ampleCone_le_nefCone}
\leanok
\uses{def:native-numerical-ample-cone,thm:native-numerical-ample-classes}
Every actual geometric ample line class belongs to the original ample
cone, and this entire cone is contained in the genuine original nef cone.
\end{lemma}
\begin{proof}
Every generator belongs to its hull. Actual ample-line positivity gives
nef membership for every generator. Since the genuine nef cone is a
convex cone, the minimality of the cone hull proves the containment.
\end{proof}

\begin{theorem}[Genuine ample openness in every original finite coefficient family {\cite[\S1.2]{BJ22}}]
\label{thm:native-numerical-finite-ample-openness}
\lean{EnvelopeProperty.NativeRealNumericalClasses.exists_ample_lineFamily_radius, EnvelopeProperty.NativeRealNumericalClasses.exists_open_ample_lineFamily_neighborhood}
\leanok
\uses{def:native-numerical-ample-cone,lem:native-numerical-ample-cone-nef,thm:native-numerical-ample-domination}
Let $X$ be integral and projective, $A$ an actual geometrically ample
line, and $(B_i)_{i\in I}$ any finite family of original invertible
lines. There is $\varepsilon>0$ such that
\[
 |a_i|<\varepsilon\quad(i\in I)
 \quad\Longrightarrow\quad
 [A]+\sum_i a_i[B_i]\in\operatorname{Amp}(X).
\]
Consequently there is an ordinary open neighborhood of zero in
$\mathbb R^I$ on which the entire original affine coefficient map takes
values in this same ample cone. No global finite dimensionality or
numerical domination hypothesis is supplied.
\end{theorem}
\begin{proof}
Actual uniform ample domination supplies a positive integer $m$ such
that $m[A]+[B_i]$ and $m[A]-[B_i]$ are classes of actual ample lines for
every $i$. Put $n=|I|$, $c=1/(n+1)$,
$\delta=c/(2m)$ and $\varepsilon=2\delta$. All three scalars are
strictly positive. If $|a_i|<2\delta$, then
$\delta+a_i/2$ and $\delta-a_i/2$ are strictly positive. The exact module
identity is
\[
 [A]+\sum_i a_i[B_i]
 = c[A]+\sum_i\bigl((\delta+a_i/2)(m[A]+[B_i])
                +(\delta-a_i/2)(m[A]-[B_i])\bigr).
\]
Indeed each summand equals $2\delta m[A]+a_i[B_i]$, and
$c+n(2\delta m)=1$. Every term on the right belongs to the original
ample cone, and the base term $c[A]$ handles the empty family as well.
Finite addition proves the assertion. Finally the set defined by the
finitely many strict inequalities $|a_i|<\varepsilon$ is open in the
ordinary finite product topology and contains zero. This proves the
neighborhood statement on the whole actual coefficient map.
\end{proof}
